\subsection{Алгоритъм за преобразуване на \textit{SVG} изображение на обобщена мрежа в \textit{TeX} формат}

\hspace*{2.5em}В този раздел се представя алгоритъм за автоматично преобразуване на \textit{SVG} изображение на ОМ в \textit{TeX} код, с което се решава~\ref{task:svg-to-tex}.

\subsubsection{Въведение и цел на алгоритъма}
\hspace*{2.5em}\textit{Scalable Vector Graphics (SVG)} е популярен векторен формат за двумерни изображения, който се мащабира без загуба на качество и е удобен за генериране и редакция \cite{SVG_F}. В академичните текстове обаче най-често се работи с \textit{TeX}, където фигурите трябва да се вписват чисто в документа и да изглеждат консистентно със шрифтовете и стила му \cite{Knuth1986TeXProgram}. Затова е необходим метод, с който пренасянето на \textit{SVG} изображение на ОМ в \textit{TeX} документ да е без ръчно пречертаване.

Алгоритъмът извлича позициите, преходите и връзките от входния \textit{SVG} и създава еквивалентна чрез примитиви като \texttt{\textbackslash put}, \texttt{\textbackslash circle}, \texttt{\textbackslash line} и \texttt{\textbackslash vector}. Полученият код може директно да се вмъкне в научен документ, като се избягват ръчни корекции.

В следващите раздели се описват използваният \textit{SVG} формат и предпоставките за обработка, извличането и преобразуване на координатите, изрязване на ребрата до границите на елементите и създаването на крайния \textit{TeX} резултат.


\subsubsection{Описание на формата на входния \textit{SVG} файл}
\hspace*{2.5em}Входният \textit{SVG} файл, с който работи алгоритъмът, следва предварително зададена структура. Всеки графичен елемент съдържа атрибут \texttt{class}, по който се определя неговият тип (позиция, преход или ребро). Стойностите на \texttt{class} могат да се генерират директно от алгоритъма за изчертаване на ОМ, описан в \textbf{Раздел \ref{sec:auto-drawing}}, и се използват като основен ориентир при обработката.

В контекста на настоящата глава, както и в имплементацията на алгоритъма в новия продукт, компонентите на ОМ са описани в \textit{SVG} файла като отделни групи \texttt{<g>} в следния формат:

\begin{itemize}
	\item \textbf{\textit{Places} (позиции)}: Представени са като групи с атрибут \texttt{class} със стойност по шаблон \texttt{node place \{id\}}, където \texttt{id} е идентификаторът на позицията. В групата присъства елемент от тип \textit{ellipse} (\texttt{<ellipse>}) за визуализацията, както и елементи \textit{title} (\texttt{<title>}) и \textit{text} (\texttt{<text>}), в които е записан идентификаторът.
	
	\item \textbf{\textit{Transitions} (преходи)}: Представени са като групи с атрибут \texttt{class} със стойност по шаблон \texttt{node transition \{id\}}, където \texttt{id} е идентификаторът на прехода. Групата съдържа два елемента от тип \textit{polygon} (\texttt{<polygon>}): единият описва вертикалната черта на прехода, а вторият – триъгълника над нея; и двата са маркирани с \texttt{class="triangle"}. Идентификаторът се съдържа в \textit{title} и \textit{text}.
	
	\item \textbf{\textit{Edges} (ребра)}: Представени са като групи с атрибут \texttt{class} със стойност по шаблон \texttt{edge \{originId\}\_\_\{destinationId\}}, където \texttt{originId} и \texttt{destinationId} са идентификаторите на свързаните елементи. Групата съдържа \textit{path} (\texttt{<path>}) за точките, по които се изчерава реброто, и \textit{polygon} (\texttt{<polygon>}) за върха на стрелката.
\end{itemize}


\subsubsection{Извличане на координатите от входния \textit{SVG} файл}
\hspace*{2.5em}Алгоритъмът \textbf{разчита} само координатите, нужни за построяване на фигурата в \textit{TeX}:

\begin{itemize}
  \item \textbf{\textit{Places} (позиции)}: от \texttt{<ellipse>} се взимат \texttt{cx, cy} (център), а от \texttt{<text>} -- \texttt{x, y} (позиция на етикета).

  \item \textbf{\textit{Transitions} (преходи)}: разчитат се \textbf{и двата} \texttt{<polygon>} елемента (за вертикалната черта и триъгълника), както и \texttt{x, y} на етикета от \texttt{<text>}.

  \item \textbf{\textit{Edges} (ребра)}: от \texttt{<path>} се взима атрибутът \texttt{d}, който е \textbf{низ със стойности} за чертане: \texttt{M x y} (\textit{move to}) задава начална точка, а \texttt{L x y} (\textit{line to}) задава следваща точка, до която се изчертава права. Така \texttt{d} се преобразува в последователност от точки \texttt{(x\_k,y\_k)}. Върхът на стрелката се разчита от \texttt{<polygon points="..."{}>} като точки \texttt{(x\_j,y\_j)}.
\end{itemize}

Накрая се прочита и \textbf{глобалната трансформация} (например \texttt{translate} и \texttt{scale}), която се прилага върху всички координати и е описана по-долу.
\subsubsection{Глобални трансформации и нормализация на координатите}
\hspace*{2.5em}Суровите координати, прочетени от отделните елементи, \textbf{не са достатъчни}, защото \textit{SVG} често съдържа \textbf{обща трансформация}, която се прилага върху всички елементи наведнъж. Това означава, че реалната позиция и размер на всяка точка в изображението се получават \textbf{след} глобално мащабиране и изместване. Ако тази стъпка се пропусне, резултатът в \textit{TeX} ще бъде с грешен мащаб и разположение, а връзките няма да съвпадат коректно с позициите и преходите.

Типичната форма на глобалната трансформация е:

\begin{center}
	\texttt{transform="translate(tx, ty) scale(s)"}
\end{center}

Извличат се:
\[
\mathit{globalTranslate}_x = tx, \quad
\mathit{globalTranslate}_y = ty, \quad
\mathit{globalScale} = s
\]

За всяка точка \((x,\,y)\) в \textit{SVG} се прилага преобразуването в \eqref{eq:svg-transform}:

\begin{equation}
	\widehat{X} = x \cdot s + tx, \quad
	\widehat{Y} = y \cdot s + ty
	\label{eq:svg-transform}
\end{equation}

След това цялото изображение се намалява чрез мащабен множител, което е показано в \eqref{eq:svg-scale}. Причината е практическа: така фигурата става удобна за директно вмъкване в документ, без да се налагат допълнителни ръчни корекции на мащаба. Понеже множителят се прилага еднакво върху всички точки, пропорциите и топологията на мрежата се запазват:

\begin{equation}
	X' = \widehat{X} \cdot \underbrace{\tfrac{6}{10}}_{\text{LENGTH\_SCALE\_FACTOR}},\quad
	Y' = \widehat{Y} \cdot \underbrace{\tfrac{6}{10}}_{\text{LENGTH\_SCALE\_FACTOR}}
	\label{eq:svg-scale}
\end{equation}

След глобалните преобразувания се прави \textbf{нормализация} към координатната система на средата \texttt{picture} в \textit{TeX}. \texttt{picture} е вградена среда за елементарно векторно чертане, в която обектите се поставят с координати чрез \texttt{\textbackslash put(x,y)\{...\}}, а платното има фиксирани размери \texttt{\textbackslash begin\{picture\}(W,H)}.

Нормализацията е необходима поради следните причини:
\begin{itemize}[noitemsep, topsep=0pt]
  \item (1) В \textit{TeX} координатният произход е в долния ляв ъгъл, докато при \textit{SVG} оста \(y\) расте надолу.
  \item (2) \textit{TeX} работи най-удобно с координати, „събрани“ в рамките на платното (без отрицателни стойности).
\end{itemize}

Първо се изчислява \textbf{ограничаващ правоъгълник} (\textit{bounding box}) на всички значими точки след \eqref{eq:svg-scale}:
\[
X_{\min},\,X_{\max},\,Y_{\min},\,Y_{\max}.
\]
След това всяка точка \((X',Y')\) се преобразува към \textit{TeX} координати, както е показано в формула \eqref{eq:tex-normalization}, като едновременно се премества спрямо минималните стойности и се обръща оста \(y\):

\begin{equation}
x_{\textit{TeX}} = (X' - X_{\min})\cdot c,\qquad
y_{\textit{TeX}} = (Y_{\max} - Y')\cdot c
\label{eq:tex-normalization}
\end{equation}

Накрая размерите на \textit{\texttt{picture}} платното се задават спрямо обхвата, което гарантира, че цялата фигура се побира коректно. Конкретно, \(W\) и \(H\) се изчисляват според \eqref{eq:picture-dimensions}.

\begin{equation}
W = (X_{\max}-X_{\min})\cdot c,\qquad
H = (Y_{\max}-Y_{\min})\cdot c
\label{eq:picture-dimensions}
\end{equation}

\subsubsection{Нормализация на размери и етикети}
\hspace*{2.5em}След преобразуването на координатите към системата на \texttt{picture} средата, е необходимо да се \textbf{уеднаквят размерите} на базовите елементи и да се \textbf{коригират позициите на етикетите}. Причината е, че входният \textit{SVG} често е създаден спрямо шрифтове и мащаб в \textit{Graphviz}, докато в \textit{TeX} реалните размери на текста и елементите са различни. Ако директно се пренесат размерите, етикетите могат да се застъпват с фигурите (особено при преходите), а различни изображения да изглеждат непоследователно.

\paragraph{Фиксиране на радиуса на позициите.}
\mbox{}\\
За да има еднакъв визуален стил във всички генерирани схеми, радиусът на всяка позиция се фиксира до константа в \textit{TeX} координати. В имплементацията това означава, че независимо от първоначалния размер в \textit{SVG}, в резултата се използва фиксиран диаметър на окръжността:
\begin{equation}
\text{\texttt{\textbackslash circle\{1.0\}}} \quad \Rightarrow \quad r_{\textit{TeX}} = 0.5.
\label{eq:fixed-place-radius}
\end{equation}
Така радиусът е постоянен (вж. \eqref{eq:fixed-place-radius}), което улеснява и последващото изрязване на ребрата до границата на кръга (описано в следващите стъпки на алгоритъма).

\paragraph{Корекция на етикетите на позициите.}
\mbox{}\\
Етикетът на позиция в \textit{SVG} обикновено е поставен близо до центъра, но при прехвърляне в \textit{TeX} визуалното разстояние може да се промени. Затова етикетът се премества леко към окръжността по посоката от центъра към оригиналната му позиция. Ако \(C=(x_c,y_c)\) е центърът на позицията, а \(L=(x_l,y_l)\) е прочетената позиция на етикета, се използва:
\begin{equation}
L_{\textit{new}} = C + k\,(L-C),
\label{eq:place-label-shift}
\end{equation}
където $k$ е коефициентът \texttt{placeLabelOffsetFactor}, който е входен за алгоритъма и може да се експериментира с него. Практически \eqref{eq:place-label-shift} приближава надписа към фигурата и намалява вероятността от припокриване с други елементи.

\paragraph{Корекция на етикетите на преходите.}
\mbox{}\\
Преходите съдържат триъгълник и вертикална линия, а етикетът често попада твърде близо до триъгълника след преобразуване. За да се избегне визуален конфликт, етикетът на преход се измества нагоре с константно вертикално отместване (вж. \eqref{eq:transition-label-offset}).
\begin{equation}
y_{\textit{new}} = y_{\textit{old}} - \texttt{transitionLabelVerticalOffset}.
\label{eq:transition-label-offset}
\end{equation}
Параметърът \texttt{transitionLabelVerticalOffset} също е входен за алгоритъма и може да се експериментира, за да осигурява четимост без ръчни корекции върху резултатния файл.

В резултат от тези две настройки -- фиксиран радиус на позициите (вж. \eqref{eq:fixed-place-radius}) и контролирани отмествания на етикетите за позиции и преходи (вж. \eqref{eq:place-label-shift} и \eqref{eq:transition-label-offset}), резултатният \textit{TeX} код дава \textbf{четима} визуализация на тези елементи.

\subsubsection{Изрязване на ребрата до границите на елементите}
\hspace*{2.5em}След нормализацията на координатите и фиксирането на размерите, остава ключов практически проблем: ребрата, прочетени от \texttt{<path>}, често започват и завършват \textbf{вътре} в съответния елемент (позиция или преход). Ако тези крайни точки се пренесат директно в \textit{TeX}, линиите ще навлизат в окръжността на позицията или ще пресичат графиката на прехода, което влошава четимостта. Затова алгоритъмът \textbf{изрязва крайните точки} на всяко ребро така, че те да попаднат точно върху \textbf{границата} на свързаните елементи.

Реброто се представя като последователност от точки
\((p_0, p_1, \dots, p_{n-1})\), получени от стойностите \texttt{M} и \texttt{L} в атрибута \texttt{d}.
Корекцията се прави \textbf{само} за двата края:
\begin{itemize}
  \item началната точка \(p_0\) се премества до границата на \textbf{изходния} елемент, като се използва посоката към \(p_1\);
  \item крайната точка \(p_{n-1}\) се премества до границата на \textbf{входния} елемент, като се използва посоката от \(p_{n-2}\) към \(p_{n-1}\).
\end{itemize}

\paragraph{Изрязване на крайни точки на ребра при позиции.}
\mbox{}\\
Позициите се чертаят като окръжности с фиксиран радиус \(r_{\textit{TeX}}\) (вж. \eqref{eq:fixed-place-radius}), което позволява прост и стабилен метод. Ако \(C=(x_c,y_c)\) е центърът на позицията, а \(P=(x_p,y_p)\) е крайна точка на реброто (преди корекцията), новата точка \(P'\) се получава чрез нормализиране по посоката \(C \rightarrow P\), както е показано във формула \eqref{eq:clip-circle}:
\begin{equation}
P' = C + r \cdot \frac{(P - C)}{\lVert P - C \rVert}.
\label{eq:clip-circle}
\end{equation}
Така реброто приключва точно на окръжността, независимо откъде идва.

\paragraph{Изрязване на крайните точки на реброто при преходи.}
\mbox{}\\
Преходите са описани в \textit{SVG} чрез многоъгълници (\textit{polygons}), но за целите на изрязването се коригира \textbf{началната или крайната точка на реброто}, така че да попадне върху границата на прехода, приближена чрез ограничаващ правоъгълник
\([x_{\min},x_{\max}] \times [y_{\min},y_{\max}]\).
Тук има значение дали преходът е \textbf{в началото} или \textbf{в края} на реброто, защото посоката се определя от различен съседен сегмент: при изрязване в началото се използват \(P=p_0\) и \(Q=p_1\), а при изрязване в края: \(P=p_{n-1}\) и \(Q=p_{n-2}\).

Нека \(P=(x_p,y_p)\) е точката, която ще се коригира, а \(Q=(x_q,y_q)\) е съседната ѝ точка по реброто. Разликите по координати се дефинират с формула \eqref{eq:deltas}:
\begin{equation}
\Delta x = x_q - x_p,\qquad \Delta y = y_q - y_p.
\label{eq:deltas}
\end{equation}

За да се избере правилната страна на правоъгълника, се използва сравнението в \eqref{eq:dominant-dir}. Ако \(|\Delta x|\) е по-голямо или равно на \(|\Delta y|\), сегментът се приема за \textbf{хоризонтален} (движението е основно наляво/надясно), а иначе -- за \textbf{вертикален} (движението е основно нагоре/надолу):
\begin{equation}
\text{хоризонтален, ако } |\Delta x| \ge |\Delta y|,\qquad
\text{вертикален, ако } |\Delta y| > |\Delta x|.
\label{eq:dominant-dir}
\end{equation}

След като се определи посоката, точката се поставя на съответната страна на ограничителния правоъгълник. Изборът на страна зависи от знака на \(\Delta x\) или \(\Delta y\), както е обобщено във формула \eqref{eq:rect-clip}:
\begin{equation}
P'=
\begin{cases}
(x_{\max},\,y_p), & \text{ако } |\Delta x|\ge|\Delta y| \text{ и } \Delta x>0,\\
(x_{\min},\,y_p), & \text{ако } |\Delta x|\ge|\Delta y| \text{ и } \Delta x<0,\\
(x_p,\,y_{\max}), & \text{ако } |\Delta y|>|\Delta x| \text{ и } \Delta y>0,\\
(x_p,\,y_{\min}), & \text{ако } |\Delta y|>|\Delta x| \text{ и } \Delta y<0.
\end{cases}
\label{eq:rect-clip}
\end{equation}

По този начин реброто докосва прехода естествено (отляво/отдясно или отгоре/отдолу), без да преминава през него.

\paragraph{Запазване на коректния край на стрелката.}
\mbox{}\\
След като крайната точка на реброто се премести до границата на входния елемент, последният сегмент се генерира като \texttt{\textbackslash vector}, така че върхът на стрелката да съвпада с новата крайна точка. В резултат ребрата не навлизат в елементите и визуалната посока се запазва коректно в \textit{TeX} фигурата.

\subsubsection{Псевдокод на алгоритъма}

\hspace*{2.5em}В по-общ вид стъпките могат да се синтезират, както следва, като дефинираните константи са с примерни стойности, получени при успешно прилагане на алгоритъма върху реални изображения:
\begin{verbatim}
function svgToTeX(svgString):
  1) DOM = parse SVG string into an XML DOM object
  2) mainG = find <g id="graph0"> and parse 
  	 transform="translate(tx, ty) scale(s)"
       globalTranslate = (tx, ty)
       globalScale = s
  3) Define LENGTH_SCALE_FACTOR = 0.6   # (2/3) shrink
  4) Define pxToTeX = 0.1              # 1px -> 0.1 TeX units

  # 5) Identify PLACES
  for each g with class="node place":
    read ellipse(cx, cy)
    apply global transform + LENGTH_SCALE_FACTOR => (x', y')
    store nodeShapes[placeName] = circle with center=(x', y') 
    and forced radius=5px

  # 6) Identify TRANSITIONS
  for each g with class="node transition":
    read polygon points => bounding box => xMin, xMax, yMin, yMax
    apply transform => store nodeShapes[transitionName] = rect

  # 7) Capture LABELS for places and transitions
  for each place:
    read original label coords => (lx, ly)
    apply transform => (lx', ly')
    offset label => (lx'', ly'') = (cx', cy') 
    				+ alpha*(lx'-cx', ly'-cy')
  for each transition:
    shift label upward => labelY' = labelY - someVerticalOffset

  # 8) Identify EDGES
  for each g with class="edge":
    read origin, destination from class or <title> text
    parse path (M x0 y0 L x1 y1 ...)
    apply transform => rawPoints
    # arrow polygon => transform as well

  # 9) CLIP edge endpoints
  for first and last rawPoints:
    if shapes[origin] == circle => move boundary
    if shapes[destination] == rect => clamp to bounding box
    (similarly for the other end)

  # 10) Compute bounding box over all x' and y'
  xMinAll, xMaxAll, yMinAll, yMaxAll

  # 11) Convert to TeX coords => (x_T, y_T)
     x_T = (x' - xMinAll)*pxToTeX
     y_T = (yMaxAll - y')*pxToTeX

  # 12) Build a \begin{picture}(width, height) with lines for edges,
      circles for places, etc.

  return final TeX code as a string
\end{verbatim}


След изпълнението на тези стъпки се генерира \textit{TeX} фрагмент, който може директно да се вмъкне в научен документ, в съответствие с приетата конвенция за изчертаване на ОМ.

\subsubsection{Пример за резултат от алгоритъма в \textit{TeX} формат}
\hspace*{2.5em}След приключване на всички преобразувания резултатът най-често се поставя във вътрешността на \texttt{\textbackslash begin\{picture\}} ... \texttt{\textbackslash end\{picture\}} блок. Примерен фрагмент:
\begin{verbatim}
\tiny
\begin{center}
\setlength{\unitlength}{8pt}
\begin{picture}(width,height)
  \linethickness{0.4pt}

  % --- Places ---
  \put(10.00,12.00){\makebox(0,0)[cc]{p1}}
  \put(10.00,12.00){\circle{1.0}}

  % --- Transition ---
  \put(15.00,20.00){\makebox(0,0)[b]{Z1}}
  \put(15.00,25.00){\makebox(0,0)[bb]{\Large$\bigtriangledown$}}
  \put(15.00,25.00){\line(0,-1){10.00}}

  % --- Edge ---
  % p1 -> Z1
  \put(11.00,12.00){\line(1,0){4.00}}
  \put(15.00,12.00){\vector(1,0){1.00}}

\end{picture}
\end{center}
\normalsize
\end{verbatim}

\noindent Тук:
\begin{itemize}
  \item \texttt{\string\put(x,\,y)\{...\}} задава мястото (в \textit{TeX} мерни единици, които зависят от \texttt{\string\unitlength}), където да се постави даден обект.
  \item \texttt{\textbackslash circle\{1.0\}} чертае окръжност с диаметър 1.0 (което съответства на радиус 0.5).
  \item \texttt{\textbackslash line(0,-1)\{10.00\}} рисува линия с дължина 10 (по вертикала, надолу).
  \item \texttt{\textbackslash vector(1,0)\{...\}} чертае линия със стрелка, която се използва за последния сегмент от реброто.
  \item \texttt{\{p1\}} поставя текст, центриран хоризонтално и вертикално.

\end{itemize}